\section{Commit each insurance event as a complete book transition}
\label{sec:insurance-atomic}

Separate group measurements still have to produce one dated company balance
sheet. In the preceding example, R's loss-recovery component depends on B's
loss component. Publishing B's financing while retaining R's opening
measurement would combine amounts from different stages of the calculation.
The proposed close operation therefore declares A, B and R as the required
group set for each event in this example. Each group retains its own identity
and measurement; the event publishes their consequences together in one
entity, book and branch. This is a declared dependency for these facts, not a
requirement to update every group in every insurance transaction.

The year end contains four events, in order: finance, service, claim and
recovery incurrence, then settlement. Each has its own business-event
identifier and immutable opening state. Keeping those boundaries preserves
the incurred balances and the distinction between earning income and moving
cash. The following intermediate states follow from the preceding journals.
Coverage columns show liabilities for A and B and an asset for R.
\begin{center}
\begin{tabular}{lrrrrr}
\toprule
After event & Cash & A coverage & B coverage & R coverage & Equity\\
\midrule
Opening & 250 & 120 & 130 & 85 & 85\\
Finance & 250 & 131 & 142 & 93 & 70\\
Service & 250 & 65.5 & 71 & 46.5 & 160\\
Incurrence & 250 & 65.5 & 71 & 46.5 & 83\\
Settlement & 173 & 65.5 & 71 & 46.5 & 83\\
\bottomrule
\end{tabular}
\end{center}
Service releases the coverage amounts before the separate claim expenses
and recovery income are recognised. Incurrence creates liabilities of 44
and 66 and a recovery asset of 33. Those balances explain why equity is 83
at that stage even though cash remains 250. Settlement clears them and
moves cash by \(-44-66+33=-77\), leaving equity unchanged. Collapsing the
four events into a single final balance would conceal this distinction.

\subsection*{A valid component calculation is part of permission to post}

The proposed operation uses the common
\texttt{OperationRequest} for \texttt{accounting.commit@1}. Its four arguments
remain references to the accounting basis, posting authorization, journal
proposal and economic delta. The expected company head remains in the
common request envelope. Each group measurement identifies the group,
issued or held role, covered group where applicable, event phase, opening
and proposed closing components, book, branch, time, information view,
method and source. A complete measurement has actual component values;
an unresolved one has no closing value and gives its reason.

Before looking up a previous receipt, the operation validates the complete
group set, common ownership and information view, admitted evidence,
measurement calculation and posting authorization. If R is unresolved,
the whole declared event lacks permission. A and B's completed proposals
remain available for diagnosis, but they do not form a closing company
state. A missing group, duplicate group, wrong branch or different method
also prevents lookup. The existing receipt index cannot turn malformed
historical input into an eligible request. This proposed profile admits
only the authored synthetic case; submitting it to a real book fails that
permission check.

For an eligible new event, the journal and economic delta must then agree
with the measurements. During finance, R's asset rises by 8 because its
claim present value gains 6 and its cost CSM gains 2. Its loss-recovery
component separately rises by \(18/13\), as a component within the asset.
An extra debit to R's asset and credit to finance income for \(18/13\)
would balance, yet overstate both the asset and income. The operation
rejects that proposal. It also rejects a proposal whose journal is correct
but whose economic delta leaves the recovery component at its old value.
The monetary accounts and component memory must both be correct.

Publication inserts the journal, complete component and account changes,
closing company state, their dependency records and edges, and the immutable
receipt, then advances the company head in the same transaction. The finite
exchange examples list those published records separately from the global
registry of proposed records. Merely retaining a rejected measurement for
inspection therefore does not make it part of the book. A rejected request
publishes none of its proposed transition. This requirement has one legal
book as its owner; it supplies no transaction guarantee across different
books or entities.

\subsection*{Retry the event without repeating its accounting}

The business-event lookup key remains the entity, book, branch and
business-event identifier. After input validation, an equal request
fingerprint returns the original receipt before testing the freshness of
today's head. Suppose all four events have committed, advancing the head
from generation zero to generation four. Retrying the original finance
request returns its original receipt and closing state at generation one.
The response separately reports the observed generation-four head and that
this invocation made no additional commit. It does not apply finance again
or replace today's state with the old closing state.

Equality covers the complete request, including its request identifier,
budget, original expected head and dependency identities. Changing any of
those bytes while retaining the business-event key is a conflict, not an
exact retry. A previously unseen event with an old head is stale. The
transaction repeats both the key and head checks before publication: an
identical request that has just won returns its receipt; different request
bytes for that event conflict; a different event that advances the head
makes this request stale. In each case the competing transaction's changes
are distinguished from any changes made by the requester.

If the acknowledgement disappears, the client cannot infer whether the
transaction committed. The exchange therefore records a transport
observation with unknown commit status, the original submission and its
exact request bytes. It does not invent a received operation result. One
possible storage state still has the opening head and no receipt; the other
has the complete closing state and receipt. Resubmitting the exact request
commits once in the first case and returns the old receipt in the second.
A timeout supplies no reason to change the event identifier, opening head
or economic facts.

These finite examples check the equations, complete record bindings and
required decisions for both storage states. They specify observable
behavior; actual transaction isolation, crash recovery and durability need
later implementation tests. Their completed coverage is a decision on one
request, with no interpretation as probability mass over financial histories.
The underlying source applicability and supplied actuarial assumptions of
the connected-group example continue to govern its accounting meaning.
